% §4.4.4 Case 3d: Vendor confirms mainnet applicability
% Swap this file in at Week 23 if vendor provides substantive technical confirmation.
% Probability: 7.5% (v8.A-final outcome distribution)

\subsubsection{Coordinated Disclosure and Vendor Response}
\label{sec:pefo-disclosure-3d}

We initiated coordinated disclosure with \DEPLOYEDAPP{} on $D_3$
(\S\ref{sec:disclosure}). The disclosure package included: (i) the structural
condition analysis in \S\ref{sec:pefo-structural}, (ii) the testnet exploit
in \S\ref{sec:pefo-testnet}, and (iii) the mainnet applicability analysis
in \S\ref{sec:pefo-mainnet}.

\paragraph{Vendor Response.}
The \DEPLOYEDAPP{} team acknowledged the vulnerability within [X] days and
confirmed that the mainnet architecture is subject to the PEFO condition
identified in \S\ref{sec:pefo-mainnet}. Specifically, the team confirmed:

\begin{enumerate}
  \item [SPECIFIC ARCHITECTURAL CONFIRMATION — fill from actual response]
  \item [SPECIFIC MAINNET DISTINGUISHING FACT CONFIRMED — fill from actual response]
  \item [MITIGATION STATUS — fill from actual response]
\end{enumerate}

The team proposed [MITIGATION / TIMELINE — fill from actual response].
We agreed to coordinate public disclosure at [DATE — fill from actual response],
contingent on [CONDITIONS — fill from actual response].

\paragraph{Disclosure Timeline.}
Table~\ref{tab:disclosure-timeline-3d} summarizes the disclosure timeline.
All events are independently verifiable via the OpenTimestamps proofs in
Appendix~\ref{app:disclosure}.

\begin{table}[h]
\centering
\caption{Disclosure timeline (\DEPLOYEDAPP{}, Case 3d).}
\label{tab:disclosure-timeline-3d}
\begin{tabular}{ll}
\toprule
Event & Date \\
\midrule
Initial outreach ($D_1$) & [DATE] \\
Formal disclosure ($D_3$) & [DATE] \\
Vendor acknowledgment & [DATE] \\
Mainnet confirmation & [DATE] \\
Agreed public disclosure & [DATE] \\
\bottomrule
\end{tabular}
\end{table}

\paragraph{Implication.}
The vendor's mainnet confirmation independently validates the architectural
analysis in \S\ref{sec:pefo-mainnet}. We note that this confirmation does not
affect the logical validity of Theorem~\ref{thm:a2-necessary}; the structural
PEFO condition follows from the architecture, not from vendor acknowledgment.
The confirmation provides an additional evidence type (vendor-authenticated
architecture) that strengthens the empirical case.
